\documentclass[conference]{IEEEtran}
\IEEEoverridecommandlockouts
\usepackage{cite}
\usepackage{amsmath,amssymb,amsfonts}
\usepackage{graphicx}
\usepackage{booktabs}
\usepackage{url}
\usepackage[hidelinks]{hyperref}

\begin{document}

\title{Measuring Eval-Set Self-Poisoning: A Synthetic Benchmark for Provenance Taint in
On-Device Verifier Pipelines}

\author{\IEEEauthorblockN{Agni Tripathi}
\IEEEauthorblockA{Independent Researcher\\
agnitripathi2022@kelloggalumni.northwestern.edu}}

\maketitle

\begin{abstract}
On-device agentic systems increasingly gate their own actions with a local verifier that records
a reference answer and later grades against it. We study a failure that lives inside a single such
pipeline: a verifier writes its own model-mediated inference into an evaluation log, that value is
promoted to ``gold'' without provenance discipline, and a downstream gold-dependent decision then
consumes it as ground truth. We call this eval-set self-poisoning through provenance taint. We
build a fully synthetic, self-contained benchmark that induces the failure under controlled
conditions and measures it with a real large language model verifier whose decisions are scored
against an immutable oracle held outside the writable pipeline. On matched tainted and clean
episodes with explicit provenance identifiers, we report a taint-adoption rate (how often the
grader follows the self-authored value against the oracle), and we ablate a guardrail into its two
parts, a provenance tag and the authoritative source, to see which drives resistance. Across four
models from two providers (pooled $n=240$ each) the baseline grader adopts the self-authored value
in almost all eligible episodes. Surfacing the source removes adoption for every model. A provenance
tag on its own is model-dependent: it nearly eliminates adoption for one model yet barely moves a
small model that keeps trusting its own unverified gold. All results are bounded to the synthetic
pipeline and the models tested; we make no claim about real-world prevalence. The benchmark and
harness are released as a reproducible artifact.
\end{abstract}

\begin{IEEEkeywords}
LLM-as-judge, verifier reliability, evaluation integrity, provenance, on-device agents,
trustworthy AI
\end{IEEEkeywords}

\section{Introduction}
On-device agentic AI runs its evaluation loop close to the action. A local verifier reads context,
records a reference answer for a task, and later grades a candidate against that record. The record
is convenient. It is also dangerous, because the verifier authored it.

We study a specific failure inside one pipeline pass. In a first step the verifier records a
reference (``gold'') answer from provided material, as a retrieval-grounded verifier would. When
that material is imperfect, the recorded value is wrong. With no discipline about where the value
came from, a provisional, model-inferred reading is promoted to ground truth. In a second step a
downstream gold-dependent decision grades a candidate against the stored value and, if it defers to
the record, marks a correct candidate wrong. The pipeline is now deferring to a value the pipeline
itself wrote. We call this eval-set self-poisoning through provenance taint.

Model collapse describes a training-scale process, and judge leakage involves two related models.
The failure here needs only one verifier inside one pass: author, store, consume. The question we
ask is empirical. Given a real verifier and an imperfect logging step, how often does a
self-authored value get adopted as ground truth, and does a lightweight provenance guardrail catch
it without wrecking clean cases?

We make three contributions.
\begin{itemize}
\item A fully synthetic, self-contained benchmark that induces provenance taint under controlled,
matched conditions with explicit provenance identifiers, and an immutable external oracle that is
architecturally held outside the writable pipeline.
\item A measurement protocol whose central number, the taint-adoption rate, is a real verifier's
decision scored against the oracle, not a scripted rate, together with the recall and false-alarm
cost of a provenance-tag plus quote-source guardrail on clean controls.
\item An empirical result across four models from two providers showing that the failure is
inducible and measurable, and an ablation showing that the authoritative source, not the provenance
label alone, is what reliably prevents adoption, with the label's effect ranging from near-total to
negligible across models. All claims are bounded to the tested synthetic pipeline and models.
\end{itemize}

\section{Related Work and Distinction}
Adjacent contamination loops are documented, but not the within-pipeline log-to-eval-key mechanism
measured here.

\textbf{Self-authored evidence and model collapse.} RAG Collapse shows responses degrade when
retrieved documents are self-authored~\cite{ragcollapse}, an analogue at the retrieval layer rather
than the evaluation key. Epidemiological modeling of model collapse studies training-scale
synthetic-data contamination~\cite{modelcollapse}. Neither concerns a verifier writing its own
inference into a gold record and consuming it within one pass.

\textbf{Judge contamination and self-preference.} Preference Leakage identifies contamination when a
judge and a data generator are related~\cite{prefleak}. That is a cross-model effect; our taint is
intra-pipeline and requires only one verifier.

\textbf{Oracle faults in synthetic judging.} The Test Oracle Problem in synthetic LLM-as-judge
corpora shows generator faults can distort or erase the intended signal~\cite{oracleproblem}, and
work auditing feedback-driven evolution shows a generated single-reference oracle can inflate
measured gains~\cite{feedbackevo}. These concern faults in how the oracle or corpus is generated. We
instead hold the oracle fixed and measure whether a self-authored value is adopted against it.

\textbf{Self-poisoning of agent artifacts.} EVOMAL demonstrates self-poisoning of a shared
coding-agent skill library, where malicious skills propagate through retrieve-and-imitate
loops~\cite{evomal}. The object there is a shared skill library. Our failure concerns an
evaluation key that a downstream decision consumes as gold.

\textbf{Construction-side avoidance.} Ground Truth First inverts benchmark construction so gold
facts are emitted before any model text, validated by a fidelity verifier, which avoids self-poisoning
by design~\cite{groundtruthfirst}. This is complementary. It prescribes how to build a clean
benchmark; it does not provide an instrument that measures the failure when a pipeline is not built
gold-first. Our benchmark is that instrument.

In short, prior work either documents adjacent loops or avoids the failure by construction. We
provide a controlled measurement of how often a real verifier's self-authored value is adopted as
ground truth against a held-out oracle, and how well a provenance guardrail catches it.

\section{Benchmark Design}

\subsection{Episodes}
Each episode is a deterministic, recomputable task with an unambiguous answer: modular arithmetic,
unit conversion at a synthetic rate, base conversion, and arithmetic-sequence terms. Values are
multi-step, which places the verifier in the realistic regime where it leans on provided reference
material rather than recomputing. Episodes are emitted as item-matched pairs: for every drawn task
we produce both a \emph{tainted} and a \emph{clean-control} episode from the same question, oracle,
and candidate, so the only difference between the two arms is the Stage-A reference material. In the
tainted arm that material contains a distractor asserting a specific wrong value; in the clean-control
arm it is authoritative. Because the pairing is item-for-item, each template appears equally in both
arms, and adoption and false-alarm are measured on the same task distribution. Everything is synthetic
and seeded.

\subsection{Two-stage pipeline}
\textbf{Stage A, reference logging (writable).} The verifier records a gold value from the Stage-A
material. The logged value is then grounded against the authoritative source: it is tagged
\texttt{source-grounded} with a verbatim quote only if the source states exactly that value, and
\texttt{model-inferred} otherwise. The tagged value is written to a writable store with a provenance
identifier. A verifier that logs a wrong value from imperfect context is therefore tagged
\texttt{model-inferred} with no supporting quote, which is the signal the guardrail later uses.

\textbf{Stage B, grading (the downstream decision).} The verifier grades a candidate, set to the
true answer, against the stored gold. We run a $2\times2$ ablation over what the grader is shown,
to separate whether resistance to a self-poisoned gold comes from the verification-provenance
\emph{tag} or from the authoritative \emph{source}: \texttt{baseline} sees only the stored value;
\texttt{tag\_only} adds the provenance tag with an instruction to distrust an unverified reference,
but no source; \texttt{source\_only} adds the source but no provenance tag; \texttt{guardrail} adds
both, and is instructed to trust the reference only if it is source-grounded and to flag a
model-inferred reference that conflicts with the source. The provenance tag is a binary
verification-status signal (source-grounded versus model-inferred), not a claim about who authored
the reference.

\subsection{Oracle isolation}
The oracle holds the true answer for every episode. It is built once, exposes read-only access, and
is imported only by the scorer. The verifier and the writable store never import it and are never
handed oracle values; the verifier sees only episode fields that exclude the oracle. The true answer
does appear inside prompts, as the candidate under grading and inside the authoritative source the
guardrail may consult, so the invariant we enforce is that the oracle is never presented to the
verifier labelled as gold. This is enforced architecturally and checked by tests.

\subsection{Metrics}
All metrics are bounded to the synthetic pipeline and the tested model.
\begin{itemize}
\item \textbf{Eligibility rate.} Fraction of tainted episodes whose Stage-A logged value actually
diverges from the oracle. Only eligible episodes can exhibit taint, so we report the rate rather
than assume it.
\item \textbf{Taint-adoption rate.} Fraction of eligible tainted episodes graded \texttt{incorrect},
meaning the grader followed the self-authored wrong value and rejected the true candidate. Reported
per condition. This is a real verifier decision.
\item \textbf{Prevented adoption (recall).} Fraction of eligible tainted episodes a condition did
not adopt, by grading \texttt{correct} or by an explicit \texttt{flag}. We also report the stricter
explicit-flag rate.
\item \textbf{False-alarm rate.} On clean controls with correct, source-grounded gold, the fraction
a condition wrongly rejects or flags. This is a condition's utility cost.
\end{itemize}

\section{Experimental Setup}
We run the benchmark with four models as the verifier across two providers: Claude Haiku 4.5 and
Claude Sonnet 4.6 (Anthropic), and GPT-4o-mini and GPT-4.1 (OpenAI). Each model is run at $n=80$
episodes (40 item-matched tainted/clean pairs) for each of three seeds (7, 17, 29), pooled to
$n=240$ per model, temperature $0$. Each episode uses five verifier calls: one Stage-A logging call
and four Stage-B grading calls (the four ablation conditions). The same model plays both stages,
matching the single-pipeline setting. Reference material, distractors, and the oracle are identical
across models, so any difference in the reported rates is a property of the verifier, not the data.
We report Wilson 95\% confidence intervals over the pooled eligible set.

\section{Results}

\begin{table}[t]
\caption{Taint-adoption rate per ablation condition (synthetic pipeline, pooled $n=240$ per model,
120 eligible tainted episodes). Lower is better. False-alarm on clean controls was $0.00$ for every
model under \texttt{baseline}, \texttt{tag\_only}, and \texttt{source\_only}, and $0.00$--$0.05$
under \texttt{guardrail}. Eligibility was $\geq 0.98$ for all models.}
\label{tab:main}
\centering
\begin{tabular}{lcccc}
\toprule
Condition & GPT-4o-mini & GPT-4.1 & Haiku 4.5 & Sonnet 4.6 \\
\midrule
\texttt{baseline}    & 1.00 & 1.00 & 0.93 & 0.99 \\
\texttt{tag\_only}   & 0.97 & 0.40 & 0.03 & 0.16 \\
\texttt{source\_only}& 0.06 & 0.00 & 0.05 & 0.00 \\
\texttt{guardrail}   & 0.00 & 0.00 & 0.00 & 0.00 \\
\bottomrule
\end{tabular}
\end{table}

Table~\ref{tab:main} reports taint-adoption per condition. Under \texttt{baseline} every model
adopts the self-authored value in almost all eligible episodes (0.93 to 1.00). The grader rejects a
candidate that matches the oracle, because the stored gold, which the pipeline itself wrote,
disagrees with it. The grader had the question in hand and could have recomputed; under baseline it
deferred to the record.

The ablation separates the two ingredients of the guardrail. Showing the authoritative source
(\texttt{source\_only}) drives adoption to near zero for every model (0.00 to 0.06). Showing only
the provenance tag (\texttt{tag\_only}), with no source, has a sharply model-dependent effect:
adoption falls to 0.03 for Haiku 4.5 (Wilson 95\% CI [0.01, 0.07]) but stays at 0.97 for GPT-4o-mini
([0.93, 0.99]), with GPT-4.1 (0.40, [0.32, 0.49]) and Sonnet 4.6 (0.16) in between. A weaker verifier
told that its stored gold is unverified still adopts it. The full \texttt{guardrail} (tag and
source) removes adoption for all models (0.00, upper CI $\leq 0.03$), at a false-alarm cost of at
most 0.05 on clean controls. All rates are over 119 to 120 eligible episodes per model and are
reported as existence-and-magnitude results on this synthetic pipeline, not population estimates.

\section{Discussion}
Three findings hold on this synthetic pipeline. First, provenance taint is inducible: when the
logging step reads imperfect material, the recorded gold diverges from the oracle, and a baseline
grader that defers to the record propagates that divergence into a wrong grade, across every model
tested. Second, the failure is a genuine model decision. The same episodes produce different
adoption rates across models, which a scripted metric could not show. Third, and most useful, the
two parts of the guardrail do different work. The authoritative source reliably prevents adoption
everywhere. A provenance label on its own is not a reliable fix: its effect ranges from near-total
(Haiku 4.5) to negligible (GPT-4o-mini). A weaker verifier needs the source; a warning that its
reference is unverified is not enough on its own.

The practical reading for on-device verifier pipelines, where the verifier is often a small model,
is that a stored reference should carry its provenance and the grader should be able to consult the
source it was supposed to be grounded in. A provenance tag alone may not move a small verifier. The
cost of surfacing both is a false-alarm rate, at most 0.05 here, that a deployment would tune.

\section{Limitations}
The benchmark is synthetic and small in scope. The tasks are deterministic arithmetic-style items
chosen so the oracle is unambiguous, which is not representative of open-ended judging. The taint is
induced through a distractor in the logging step, a controlled proxy for imperfect retrieved
context. Every number is bounded to the tested pipeline and models. We make no estimate of how often
provenance taint occurs in any deployed system; establishing prevalence would require instrumented
field data, which this work does not have. The guardrail we test is deliberately simple, and its
false-alarm cost should be read as an existence result about the tradeoff, not an optimized
operating point.

\section{Conclusion}
We framed eval-set self-poisoning through provenance taint as a within-pipeline failure, built a
self-contained synthetic benchmark that induces it, and measured it with a real verifier scored
against an immutable oracle held outside the writable pipeline. The taint-adoption rate and the
recall and false-alarm cost of a provenance guardrail give a first controlled instrument for a
failure that adjacent work either documents in other settings or avoids by construction. The next
step is to widen the task distribution beyond deterministic items and to study guardrails that keep
the catch rate while lowering false alarms.

\section*{Artifact and Reproducibility}
The generator, oracle, verifier harness, guardrail, and scoring are released as a self-contained
repository with a deterministic mock verifier for continuous integration and a real-provider mode.
All content is synthetic.

\begin{thebibliography}{00}
\bibitem{ragcollapse} G. Druck and E. Smith, ``RAG Collapse: LLM Responses Collapse When Retrieved
Documents Are Self-Authored,'' arXiv:2608.22118, 2026.
\bibitem{modelcollapse} X. Wang, ``Epidemiology of Model Collapse: Modeling Synthetic Data
Contamination via Bilayer SIR Dynamics,'' arXiv:2606.05168, 2026.
\bibitem{prefleak} D. Li \emph{et al.}, ``Preference Leakage: A Contamination Problem in
LLM-as-a-judge,'' in \emph{Proc. ICLR}, 2026. arXiv:2502.01534.
\bibitem{oracleproblem} S. Ball\i, ``The Test Oracle Problem in Synthetic LLM-as-Judge Corpora:
Disappearance, Distortion and a Validation Protocol,'' arXiv:2607.13707, 2026.
\bibitem{feedbackevo} Y. Liang \emph{et al.}, ``Auditing and Decomposing Feedback-Driven Evolution
in LLM Test Generation under the Oracle Problem,'' arXiv:2608.19626, 2026.
\bibitem{evomal} X. Wu \emph{et al.}, ``EVOMAL: Self-Poisoning in Self-Evolving Coding Agents,''
arXiv:2608.25776, 2026.
\bibitem{groundtruthfirst} ``Ground Truth First: A Longitudinal Evaluation Instrument for Agent
Memory,'' arXiv:2607.21962, 2026.
\end{thebibliography}

\end{document}
